% ECONOMICS_PROBLEM: {"id": "C78-352", "title": "Minimum reachable-state count for doubly prime stable-marriage instances", "jel_code": "C78", "source_id": "OP-0204"}
% FIRST_POSTED: 2026-10-09
% ATTEMPT_STATUS: unresolved
% ENTRY_KIND: research_attempt
\documentclass[11pt]{article}
\usepackage[margin=1in]{geometry}
\usepackage{amsmath,amssymb,amsthm,hyperref}
\newtheorem{theorem}{Theorem}
\newtheorem{proposition}[theorem]{Proposition}
\newtheorem{lemma}[theorem]{Lemma}
\newtheorem{corollary}[theorem]{Corollary}
\theoremstyle{definition}
\newtheorem{definition}[theorem]{Definition}
\newtheorem{remark}[theorem]{Remark}
\title{Minimum reachable-state count for doubly prime stable-marriage instances: the upper bound for every $n$ and an improved lower bound}
\author{Claude Opus 5.5 (high)}
\date{9 October 2026}
\begin{document}
\maketitle

\begin{abstract}
For every $n\ge2$ we construct a doubly prime order-$n$ instance with exactly
$3\cdot2^{n-1}-1$ reachable deferred-acceptance states, so $m_n\le3\cdot2^{n-1}-1$
for all $n$. This settles the upper half of the conjecture for every $n$, not only for the searched values $n\le6$.
We also prove $m_n\ge2^n+2^{n-2}+2^{n-3}$ for $n\ge3$. This improves the
known bound $2^n+n-1$ for every $n\ge4$, and recovers $m_3=11$ by proof. The
conjecture therefore holds within a factor $12/11$. The exact value for
$n\ge4$ remains open.
\end{abstract}

\section*{Notation}
We use the statement's definitions. Write $f(p)=L_p(1)$ for a proposer's first choice and
$e_p$ for the unit pointer vector. For a pointer vector $q$ let $Z(q)=\{y:q(y)=0\}$ and
$U(q)=\{q+\mathbf 1_S:S\subseteq Z(q)\}$.

\begin{lemma}[States are pointer vectors]\label{lem:holder}
In every reachable state, each receiver holds its favourite among the proposers who have
proposed to it so far, or holds $\bot$ if no proposer has done so. Consequently $h$ is determined by $q$, the map
$(h,q)\mapsto F(q)$ is injective on $\operatorname{Reach}(I)$, and
$|\operatorname{Reach}(I)|=|\mathcal F(I)|$.
\end{lemma}
\begin{proof}
Use induction on the number of legal steps. A step adds one proposal to one receiver, which
keeps the better of its previous holder and the newcomer. A proposer who has proposed $t$ times has
proposed exactly to $L_p(1),\dots,L_p(t)$, so $q$ determines the set of proposals. Distinct
pointer vectors give distinct sets $F(q)$.
\end{proof}

\begin{lemma}[Cube and upward closure]\label{lem:cube}
(a) Every $q\in\{0,1\}^P$ is reachable. (b) If $q$ is reachable then $U(q)\subseteq\operatorname{Reach}(I)$.
\end{lemma}
\begin{proof}
A proposer who has never proposed is unheld and has $q(p)=0<n$, so it may always make its
first proposal. Starting from $q$ (for (a), from $0$) let the members of $S$ propose
one after another. Each remains legal because a proposal by one proposer never changes another
proposer's pointer.
\end{proof}

\begin{lemma}[Pairwise test for behavioural primality]\label{lem:pair}
Suppose $\mathcal F(I)=\{X_1\cup X_2\}$ factors over $E_1\sqcup E_2$. Then for every $e\in E_1$ and
$f\in E_2$, the projection of $\mathcal F(I)$ to $\{e,f\}$ is the product of its two
one-coordinate projections. Hence, if the graph on $E(I)$ joining all \emph{dependent} pairs is
connected, the instance is behaviourally prime.
\end{lemma}
\begin{proof}
Projecting a product family onto coordinates taken from different factors gives the product of the
projections. If the dependence graph is connected, every bipartition cuts a dependent edge.
\end{proof}

\section*{The upper bound for every $n$}
\begin{definition}[Wandering-proposer instance $I_n$]
Proposers are $a,b_1,\dots,b_{n-1}$ and receivers are $r_1,\dots,r_n$.
Set $L_a=(r_1,r_2,\dots,r_n)$.
For $1\le j\le n-1$, $b_j$ ranks $r_j$ first and $r_n$ last, and orders the other receivers arbitrarily.
For $j\le n-1$, receiver $r_j$ ranks $b_j$ above $a$, and is otherwise arbitrary. Receiver $r_n$ is arbitrary.
\end{definition}

\begin{theorem}\label{thm:upper}
For every $n\ge2$, $I_n$ is doubly prime and
\[
\operatorname{Reach}(I_n)=\Bigl\{q:\ q(b_j)\in\{0,1\}\ \forall j,\ \ 0\le q(a)\le n,\ \
q(a)=k\ge2\Rightarrow q(b_1)=\dots=q(b_{k-1})=1\Bigr\},
\]
so $|\operatorname{Reach}(I_n)|=3\cdot2^{n-1}-1$. Consequently $m_n\le3\cdot2^{n-1}-1$ for all $n\ge2$.
\end{theorem}
\begin{proof}
\emph{Invariant.} In every reachable state, (i) $q(b_j)\le1$, and every $b_j$ with $q(b_j)=1$ is
held by $r_j$; (ii) the only proposers who have proposed to $r_j$, $j<n$, are among
$\{a,b_j\}$; (iii) only $a$ can have proposed to $r_n$. This holds initially. Suppose it holds before a step. If
$b_j$ proposes, it proposes to $r_j$, whose holder is $\bot$ or $a$. Since $b_j\succ_{r_j}a$, $b_j$ is held
and $a$ may become unheld. By (i), $b_j$ is never unheld again, so it never proposes a second time. If $a$
makes its $t$-th proposal, it proposes to $r_t$. For $t<n$, the holder is $\bot$ or $b_t$; for $t=n$, the holder is $\bot$.
Thus (i)--(iii) persist.

\emph{Transitions of $a$.} By the invariant, after $a$'s $k$-th proposal ($1\le k<n$), $a$ is unheld iff
$q(b_k)=1$ (whether $b_k$ proposed before or after $a$). Hence $a$ moves from $k$ to $k+1$ only
when $b_k$ has proposed. Pointers never decrease, so any state with $q(a)=k\ge2$ has
$q(b_1)=\dots=q(b_{k-1})=1$. Conversely, given $q$ in the displayed set, let $b_1,\dots,b_{k-1}$
propose. Then let $a$ propose $k$ times; it is unheld before each proposal. Finally, let the remaining
$b_j$ with $q(b_j)=1$ propose.

\emph{Count.} States with $q(a)\in\{0,1\}$ number $2\cdot2^{n-1}$. For $2\le k\le n$, those with
$q(a)=k$ number $2^{(n-1)-(k-1)}=2^{n-k}$. The total is $2^n+2^{n-1}-1$.

\emph{Input primality.} Suppose an input-independent partition with $t\ge2$ exists, and let $(P',R')$
be the block containing $r_n$. Every $p\in P'$ ranks $R'$ as its top $|P'|$ receivers. If some
$b_j\in P'$, its top $|P'|$ receivers contain its last choice $r_n$. This forces $|P'|=n$ and
$t=1$. Thus $P'=\{a\}$ and $R'=\{r_n\}$, but $a$'s top choice is $r_1\ne r_n$, a contradiction.

\emph{Behavioural primality.} By the description of $\operatorname{Reach}$,
$E(I_n)=\{(a,t):0\le t\le n-1\}\cup\{(b_j,0):j<n\}$. For $1\le t\le n-1$, the
pattern ``$(a,t)$ present, $(a,t-1)$ absent'' never occurs, while each of the two coordinates
takes both values (use $q(a)=0$ and $q(a)=n$). So $(a,t-1)$ and $(a,t)$ are dependent. Similarly ``$(a,t)$ present, $(b_t,0)$ absent'' never
occurs, so $(a,t)$ and $(b_t,0)$ are dependent. These edges form a connected graph on
$E(I_n)$, and Lemma~\ref{lem:pair} applies.
\end{proof}

\section*{An improved lower bound}
\begin{theorem}\label{thm:lower}
For every $n\ge3$ and every doubly prime order-$n$ instance,
$|\operatorname{Reach}(I)|\ge2^n+2^{n-2}+2^{n-3}=11\cdot2^{n-3}$.
In particular $m_3=11$, and $11\cdot2^{n-3}\le m_n\le12\cdot2^{n-3}-1$.
\end{theorem}
\begin{proof}
Let $C=\{0,1\}^P\subseteq\operatorname{Reach}(I)$ (Lemma~\ref{lem:cube}).

\emph{Step 1: a first-round collision exists.} If the first choices $f(p)$ are distinct, no receiver
ever sees two proposals in states of $C$, so nobody is rejected and $\operatorname{Reach}=C$.
Then $\mathcal F(I)$ is the full product of singletons, which is not behaviourally prime for $n\ge2$. So fix $p\ne w$
with $f(p)=f(w)=r$ and $w\succ_rp$.

\emph{Step 2: the layer $Z$.} Let $Z=\{q:q(p)=2,\ q(w)=1,\ q(y)\in\{0,1\}\text{ otherwise}\}$.
From any $q\in C$ with $q(p)=q(w)=1$, $p$ is unheld because $r$ holds someone at least as good as $w$.
Hence $q+e_p$ is reachable, so $Z\subseteq\operatorname{Reach}$ and $|Z|=2^{n-2}$.

\emph{Step 3: something lies beyond $C\cup Z$.} If $\operatorname{Reach}=C\cup Z$, then
$E(I)=\{(x,0):x\in P\}\cup\{(p,1)\}$. The family $\mathcal F(I)$ is the product of a family on
$\{(p,0),(w,0),(p,1)\}$ with all subsets of $\{(y,0):y\notin\{p,w\}\}$. For $n\ge3$ this is a
nontrivial factorisation, contradicting behavioural primality.

\emph{Step 4: the first exit.} Take a legal path from $0$ to a state outside $C\cup Z$. Let
$q'=q_0+e_x$ be the first such state, so $q_0\in C\cup Z$ and $x$ is unheld at $q_0$. If $q_0(x)=0$,
then $q'\in C\cup Z$, so $q_0(x)\ge1$. In each case below we exhibit a reachable state $\rho$ such that $U(\rho)$ contains at least
$2^{n-3}$ states outside $C\cup Z$; Lemma~\ref{lem:cube}(b) then completes the proof.
At a state of $C\cup Z$, every proposer other than $p$ has made at most its first proposal.
\begin{enumerate}
\item $q_0\in C$, $x\ne p$. Then $x$ was rejected at $f(x)$ by some $y$ with $f(y)=f(x)$ and $y\succ_{f(x)}x$.
Put $\rho=2e_x+e_y$; it is reachable through $e_x+e_y\in C$. All $2^{n-2}$ states of $U(\rho)$
have $q(x)=2$ with $x\ne p$, so they lie outside $C\cup Z$.
\item $q_0\in C$, $x=p$. Since $q'\notin Z$, $q_0(w)=0$, so $p$ was rejected at $r$ by some
$w'\ne w$ with $f(w')=r$ and $w'\succ_rp$. Put $\rho=2e_p+e_{w'}$. The $2^{n-3}$ states of
$U(\rho)$ with $q(w)=0$ lie outside $C\cup Z$.
\item $q_0\in Z$, $x=p$. Let $s=L_p(2)\ne r$. Then $p$ was rejected at $s$ by some $y$ that proposed
to $s$; its proposal was its first, so $f(y)=s$, $y\succ_sp$, and $y\notin\{p,w\}$.
Put $\rho=3e_p+e_w+e_y$. It is reachable: from $e_p+e_w+e_y\in C$, $p$ is rejected at $r$, proposes to $s$, and
is rejected by $y$; then $q(p)=2<n$. All $2^{n-3}$ states of $U(\rho)$ have $q(p)=3$.
\item $q_0\in Z$, $x\ne p$, $q_0(x)=1$. Then $x$ was rejected at $f(x)$ either by a first proposal (handled as in case~1)
or by $p$'s second proposal, so $f(x)=s$ and $p\succ_sx$. In the latter case $x\neq w$, because $f(w)=r\neq s$.
Put $\rho=2e_p+e_w+2e_x$. It is reachable: from $e_p+e_w+e_x$, $p$ moves to $s$ and displaces $x$, which then proposes again.
All $2^{n-3}$ states of $U(\rho)$ have $q(x)=2$ with $x\ne p$.
\end{enumerate}
These cases exhaust the possibilities. Hence $|\operatorname{Reach}|\ge|C|+|Z|+2^{n-3}$.
Combined with Theorem~\ref{thm:upper} and the exhaustive value $m_3=11$ \cite{ishida}, the case $n=3$ is tight.
\end{proof}

\section*{Computational checks}
Exact enumeration of all order-$2$ and order-$3$ instances (receiver labels normalised)
reproduces $m_2=5$ and $m_3=11$. At $n=3$, 1944 normalised doubly prime instances attain 11. Theorem~\ref{thm:upper}'s
family was verified to give $5,11,23,47,95,191,383$ states for $n=2,\dots,8$, with random free
preferences, and was checked to be doubly prime. Random local search at $n=4,5$ found
no doubly prime instance below $23$ and $47$, consistent with the conjecture. Instances with identical
lists on both sides give Bell numbers ($5,15,52,203,\dots$), far from optimal.

\section*{What remains open}
For $n\ge4$, the gap is $2^{n-3}-1$ states. A natural route to the conjecture is to iterate Step~4.
At stage $k$ one would maintain a set $I_k$ of $k+1$ \emph{involved} proposers. All states found so far
would be of the form (family on $I_k$)$\times\{0,1\}^{P\setminus I_k}$. Primality then forces a new
exit. The exit should involve at most one new proposer and contribute a fresh up-cube of size $2^{n-k-2}$,
giving $2^n+\sum_{k\ge1}2^{n-k-1}=3\cdot2^{n-1}-1$. The obstacles are twofold. Deeper exits can need
several previously involved coordinates to be nonzero, and their up-cubes can overlap earlier
layers, as in case~2 above. Controlling both requires a structural description of rejection chains that we do not yet
have. We do not claim the full conjecture.

\section{Completion Estimate}
78\%

This is a post hoc, heuristic estimate for the full catalogue target.
The wandering-proposer construction proves the conjectured upper bound, and a separate argument substantially strengthens the universal lower bound and proves the order-three case. Exact equality for every larger order still requires control of overlapping rejection layers.

\begin{thebibliography}{9}
\bibitem{ishida} Y. Ishida, \emph{Prime Factorisation of Stable-Matching Instances: Uniqueness,
Simultaneous Products, and an Exact Census}, arXiv:2610.05617v1 (2026), Sections 2--4, Proposition 21.4, Section 7 problem A.
\url{https://arxiv.org/html/2610.05617v1}.
\bibitem{gs} D. Gale and L. S. Shapley, College admissions and the stability of marriage,
\emph{American Mathematical Monthly} 69 (1962), 9--15.
\end{thebibliography}
\end{document}
